% Irodov, source PDF pp.41--42; original printed pp.32--33.
\exproblem{191}{32}
质量 $m$ 的粒子在中心势场 $U(r)=kr^2$ 中作闭合轨道运动，
$k>0$。已知粒子与中心 $O$ 的最近距离为 $r_1$，
最远点处速率为 $v_2$，求 $m$。
\exsolution
在最近点和最远点，径向速度均为零，速率完全切向。
记最远半径为 $r_2$、守恒角动量为 $L$，则
$L=mr_1v_1=mr_2v_2$。
能量守恒可写成两端同一有效势：
\[
kr_1^2+\frac{L^2}{2mr_1^2}
=kr_2^2+\frac{L^2}{2mr_2^2}.
\]
对非圆形轨道 $r_2\ne r_1$，两边移项并约去
$r_2^2-r_1^2$，得到
$L^2/(2m)=kr_1^2r_2^2$。再由
$L=mr_2v_2$ 消去 $L,r_2$：
\[
\boxed{m=\frac{2kr_1^2}{v_2^2}}.
\]
\excheck 此结论来自中心力场的角动量与总能量同时守恒；
若轨道恰为圆，$r_2=r_1$ 时不能采用上述约因子步骤，
但圆轨道向心条件同样给出此式。

\exproblem{192}{32}
长为 $\ell$ 的轻绳悬挂一小球，
初始与竖直方向夹角为 $\theta$。在这一位置使小球
获得与其所在竖直平面垂直的水平初速 $v_0$。
求使小球运动中能达到最大偏角 $90^\circ$ 的临界初速。
\exsolution
绕竖直轴的外力矩为零，轴向角动量守恒。
初始绕轴半径为 $\ell\sin\theta$，所以
$L_z=m\ell v_0\sin\theta$。
当偏角达到 $\pi/2$ 的转折点时，
沿经线方向的速度瞬时为零，只有绕竖轴的方位速率 $u$，
此时半径为 $\ell$，故
$u=v_0\sin\theta$。
以初始点为势能零点，机械能守恒：
\[
\frac12mv_0^2
=\frac12m v_0^2\sin^2\theta+mg\ell\cos\theta.
\]
由此
\[
\boxed{v_{0,\mathrm{crit}}=
\sqrt{\frac{2g\ell}{\cos\theta}}}.
\]
\excheck 需 $0\le\theta<\pi/2$。
该临界速度对应“恰好到达 $90^\circ$”的转折状态，
初速更大将越过此角度。

\exproblem{193}{32}
质量为 $m$ 的小物体被不可伸长细绳系住，
在光滑水平面上绕固定孔 $O$ 转动。
有人以恒速将细绳向下拉入孔中，故物体到孔距离
$r$ 以恒定速率减小；在 $r=r_0$ 时角速度为 $\omega_0$。
求拉线张力 $T(r)$。
\exsolution
由于细绳张力始终径向指向孔，
关于孔 $O$ 的外力矩为零，所以
\[
mr^2\dot\varphi=mr_0^2\omega_0,\qquad
\dot\varphi=\frac{r_0^2\omega_0}{r^2}.
\]
在极坐标中径向加速度为
$a_r=\ddot r-r\dot\varphi^2$。
拉入速度恒定意味着 $\ddot r=0$，
径向牛顿第二定律 $-T=ma_r$ 得
\[
\boxed{T(r)=mr\dot\varphi^2
=\frac{m\omega_0^2r_0^4}{r^3}}.
\]
\excheck 张力随半径减小迅速增大。
若拉线速度并非常数，还必须保留 $-m\ddot r$ 项。

\exproblem{194}{32}
一根不可伸长轻绳绕在半径为 $R$ 的有质量定轴滑轮上，
自由端系一小物体质量为 $m$。从静止开始释放，
求\emph{滑轮和重物组成的整个系统}相对滑轮固定轴的
总角动量随时间的变化。
\exsolution
以滑轮固定转轴为矩心，支承轴反力对该轴不产生力矩；
绳张力是系统内部相互作用，其作用力矩在滑轮与重物的
总角动量平衡中抵消。
重物始终在距轴水平距离 $R$ 的铅垂线上运动，
其重力关于轴的力矩大小为 $mgR$，方向恒定。
对整个系统应用角动量定理：
\[
\frac{\mathrm dL_{\rm total}}{\mathrm dt}=mgR,
\qquad L_{\rm total}(0)=0,
\]
因此
\[
\boxed{|L_{\rm total}(t)|=mgRt}.
\]
\excheck 结果与滑轮转动惯量无关，因为求的是\emph{整个}
滑轮—重物系统的角动量。单独求重物的角动量不能在
未知滑轮转动惯量时直接写成 $mgRt$。

\exproblem{195}{32}
均匀实心球质量 $m$、半径 $R$，从静止开始沿
倾角为 $\alpha$ 的斜面无滑动滚下。
求相对于\emph{最初接触点} $O$ 的角动量大小 $L(t)$；
若斜面完全光滑，结果如何变化？
\exsolution
对于纯滚动均匀实心球，
$I_C=2mR^2/5$，质心沿斜面加速度
$a_C=(5/7)g\sin\alpha$，
速率为 $v=a_Ct$，角速度 $\omega=v/R$。
关于初始接触点分解：
\[
\bm L_O=(\bm r_C-\bm r_O)\times m\bm v_C
+I_C\bm\omega.
\]
球心相对初始接触点的沿斜面位移与 $\bm v_C$ 平行，
故其叉积贡献为零；垂直斜面的距离始终为 $R$，
于是两部分角动量同向，大小相加为
\[
L_O=Rm v+I_C\frac vR
=\frac75mRv=mgR\sin\alpha\,t.
\]
若斜面完全光滑，球只平动不自转，
$a_C=g\sin\alpha$，则
$L_O=Rm(g\sin\alpha\,t)$，因此
\[
\boxed{L_O(t)=mgR\sin\alpha\,t}
\]
在两种情况下相同。
\excheck 本题需要区分\emph{重力力矩}与
\emph{合外力矩}：选最初接触点$O$为固定参考点，
球心沿斜面已移动$s$，则重力相对$O$
的力臂并非一直等于$R\sin\alpha$。
取沿斜面下行为$\bm e_t$、离开斜面为$\bm e_n$，
球心位置$\bm r_C=s\bm e_t+R\bm e_n$。
重力和接触支持力（$N=mg\cos\alpha$）的
关于$O$力矩中，分别有正负
$s mg\cos\alpha$项，恰好抵消；
接触静摩擦力与$\bm r_{\rm contact}=s\bm e_t$
共线，不产生关于$O$的力矩。
因此\emph{总外力矩}的大小才恒为
$mgR\sin\alpha$，对有无摩擦均成立。
不能把“总外力矩恒定”误写为“重力力矩恒定”。

\exproblem{196}{32}
一个粒子系统相对于固定点 $O$ 的总动量为 $\bm p$，
总角动量为 $\bm L_O$。另一固定点 $O'$ 的位置矢量
相对 $O$ 为 $\bm r_0$。求系统相对 $O'$ 的角动量，
并问在何种情形下该量与参考点选取无关。
\exsolution
对第 $i$ 粒子，相对 $O'$ 的位置矢量为
$\bm r_i'=\bm r_i-\bm r_0$。因此
\[
\begin{aligned}
\bm L_{O'}
&=\sum_i(\bm r_i-\bm r_0)\times\bm p_i\\
&=\sum_i\bm r_i\times\bm p_i
-\bm r_0\times\sum_i\bm p_i.
\end{aligned}
\]
故
\[
\boxed{\bm L_{O'}=\bm L_O-\bm r_0\times\bm p}.
\]
任意点的改变都不影响角动量，当且仅当
$\boxed{\bm p=\bm0}$。
\excheck 对某一特定参考点改变，若 $\bm r_0\parallel\bm p$
也可能有相同角动量；但要求\emph{任意}选点均不变则必须 $\bm p=0$。

\exproblem{197}{32}
证明柯尼希角动量分解：系统相对于固定点 $O$ 的角动量为
$\bm L_O=\widetilde{\bm L}+\bm r_C\times\bm p$，
其中 $\widetilde{\bm L}$ 是质心平动参考系中的内禀角动量，
$\bm r_C$ 是质心相对 $O$ 的位置矢量，
$\bm p$ 为总动量。
\exsolution
写出每个质点的位置和速度相对质心的分解：
\[
\bm r_i=\bm r_C+\bm\rho_i,\qquad
\bm v_i=\bm V_C+\dot{\bm\rho}_i.
\]
质心定义保证
$\sum_i m_i\bm\rho_i=\bm0$ 和
$\sum_i m_i\dot{\bm\rho}_i=\bm0$。
展开角动量总和：
\[
\begin{aligned}
\bm L_O
&=\sum_i m_i(\bm r_C+\bm\rho_i)
\times(\bm V_C+\dot{\bm\rho}_i)\\
&=M\bm r_C\times\bm V_C+
\sum_i m_i\bm\rho_i\times\dot{\bm\rho}_i,
\end{aligned}
\]
因为两个交叉项均为零。
令 $\bm p=M\bm V_C$、
$\widetilde{\bm L}=
\sum_i m_i\bm\rho_i\times\dot{\bm\rho}_i$，得
\[
\boxed{\bm L_O=\widetilde{\bm L}+\bm r_C\times\bm p}.
\]
\excheck 此分解是恒等式，并不要求系统孤立，也不要求
内禀角动量本身守恒。

\exproblem{198}{33}
质量 $m$ 的小球以速率 $v_0$ 水平正碰一个初始静止的
刚性哑铃的一端（图1.50）。哑铃两端小球质量各为 $m/2$，
端点相距 $\ell$，忽略端球尺寸，
碰撞完全弹性。求碰后哑铃\emph{质心系}角动量大小。
\exsolution
哑铃总质量为 $m$，绕自身质心转动惯量
$I=(m/2)(\ell/2)^2+(m/2)(\ell/2)^2=m\ell^2/4$。
设水平碰撞冲量为 $J$；则哑铃质心平动速度
$V=J/m$，角速度满足
$I\omega=J\ell/2$，故
$\omega=2J/(m\ell)$。
受撞端点沿碰撞方向的速度为
$V+\omega\ell/2=2J/m$；
入射小球碰后速度为 $v_0-J/m$。
完全弹性意味着两接触点分离的相对速度等于
碰前接近速率：
\[
\frac{2J}{m}-\left(v_0-\frac Jm\right)=v_0,
\]
所以 $J=(2/3)mv_0$。
哑铃质心系角动量是该冲量对质心的矩：
\[
\boxed{\widetilde L=I\omega=\frac{J\ell}{2}
=\frac13m\ell v_0}.
\]
\excheck 这里哑铃的质心平动动能和转动动能都参与了
碰撞能量分配；不能把哑铃当作固定转轴的转子。

\exproblem{199}{33}
两块相同质量 $m$ 的小圆盘位于光滑水平面上，
用初始自然长 $\ell_0$、劲度系数 $\kappa$ 的轻弹簧连接。
使其中一块沿垂直于弹簧的方向获得速率 $v_0$。
若最大相对伸长 $\varepsilon_{\max}=
(\ell_{\max}-\ell_0)/\ell_0$ 很小，
求其近似值。
\exsolution
在质心系用相对坐标分析，约化质量
$\mu=m/2$。初始相对速度大小为 $v_0$，且与
连线垂直，所以相对角动量为
$L=\mu\ell_0v_0$。
初始能量为 $\mu v_0^2/2$。
当弹簧伸长到最大时径向速度为零，设
$r=\ell_0(1+\varepsilon)$，能量守恒：
\[
\frac12\mu v_0^2
=\frac{L^2}{2\mu r^2}
+\frac12\kappa\ell_0^2\varepsilon^2.
\]
代入 $L$ 并移项：
\[
\mu v_0^2\left[1-\frac1{(1+\varepsilon)^2}\right]
=\kappa\ell_0^2\varepsilon^2.
\]
若 $\varepsilon\ll1$，
$1-(1+\varepsilon)^{-2}=2\varepsilon+O(\varepsilon^2)$，
解得非零的首阶结果
\[
\boxed{\varepsilon_{\max}\simeq
\frac{m v_0^2}{\kappa\ell_0^2}}.
\]
\excheck 右端无量纲，且须自身远小于 $1$
才符合题目的小伸长近似条件。
